\SourcePage{389}{392}
\section{Multiplet terms. Case $a$}

We now turn to the classification of molecular levels with nonzero spin $S$. In the zeroth approximation, when relativistic effects are wholly neglected, the energy of a molecule, like that of any system of particles, is independent of the spin direction (the spin is ``free''). Every energy level is therefore degenerate with respect to this direction. Relativistic effects remove the degeneracy and make the energy depend on the projection of the spin on the molecular axis. We shall refer to relativistic interactions in molecules as the \emph{spin--axis interaction}. As in atoms, its chief contribution is the interaction of the spins with the orbital motion of the electrons\footnote{Besides spin--orbit and spin--spin interactions, there is an interaction of the spin and electronic orbital motion with the rotation of the molecule. This part of the interaction is very small and is of possible interest only for terms with $S=1/2$ (see \S~81).}.

The nature and classification of molecular levels depend essentially on the relative importance of the spin--orbital-motion interaction and molecular rotation. The importance of the latter is measured by the separations between neighboring rotational levels. Two limiting cases must accordingly be considered. In one, the spin--axis interaction energy is large compared with the rotational-level differences; in the other it is small. The first is conventionally called case (or coupling case) $a$, and the second case $b$ (F.~Hund, 1933).

Case $a$ occurs most often. The main exception is provided by $\Sigma$ terms, for which case $b$ usually applies because their spin--axis interaction is weak (see below)\footnote{The normal electronic term ${}^3\Sigma$ of the $\mathrm O_2$ molecule is a special case. Its coupling is intermediate between cases $a$ and $b$ (see problem 3 in \S~84).}. For other terms, case $b$ is sometimes encountered in the lightest molecules: there the spin--axis interaction is comparatively weak and the rotational-level spacings are large because the moment of inertia is small.

Intermediate cases between $a$ and $b$ are of course also possible. Moreover, as the rotational quantum number changes, one and the same electronic state may pass continuously from case $a$ to case $b$. The spacings between neighboring rotational levels grow with the rotational quantum number. At large values they may therefore exceed the spin--axis coupling energy (case $b$), even though the lower rotational levels belonged to case $a$.

In case $a$, the classification differs little in principle from that of terms with zero spin. We first consider the electronic terms with the nuclei fixed, completely neglecting rotation. In addition to the projection $\Lambda$

\SourcePage{390}{393}
of the electronic orbital angular momentum, we must now consider the projection of the total spin on the molecular axis. This projection is denoted by $\Sigma$\footnote{It must not be confused with the symbol for terms having $\Lambda=0$.} and has the values $S,S-1,\ldots,-S$. We take $\Sigma$ to be positive when the direction of the spin projection agrees with that of the orbital angular momentum about the axis (recall that $\Lambda$ denotes the magnitude of the latter). The quantities $\Lambda$ and $\Sigma$ add to give the total electronic angular momentum about the molecular axis:
\begin{equation}\Omega=\Lambda+\Sigma.\tag{83.1}\end{equation}
It takes the values $\Lambda+S,\Lambda+S-1,\ldots,\Lambda-S$. Thus an electronic term with orbital angular momentum $\Lambda$ splits into $2S+1$ terms distinguished by their values of $\Omega$. As for atomic terms, this splitting is called the fine structure, or multiplet splitting, of the electronic levels. The value of $\Omega$ is written as a subscript on the term symbol; for example, $\Lambda=1$, $S=1/2$ gives the terms ${}^2\Pi_{1/2}$ and ${}^2\Pi_{3/2}$.

Allowing for nuclear motion produces vibrational and rotational structure for each of these terms. The rotational levels are characterized by the quantum number $J$, the total angular momentum of the molecule, including the orbital and spin angular momenta of the electrons and the rotational angular momentum of the nuclei\footnote{As usual, $K$ is reserved for the total molecular angular momentum excluding spin. In case $a$ there is no quantum number $K$, since $\mathbf K$ is not even approximately conserved.}. The number $J$ takes integral values beginning with $|\Omega|$:
\begin{equation}J\geqslant|\Omega|.\tag{83.2}\end{equation}
This is the evident generalization of (82.6).

We derive quantitative formulas for the molecular levels in case $a$. First consider the fine structure of an electronic term. In treating the fine structure of atomic terms in \S~72, we used (72.4), according to which the mean spin--orbit interaction is proportional to the projection of the atom's total spin on the orbital-angular-momentum vector. Quite analogously, the spin--axis interaction in a diatomic molecule, averaged over the electronic state at a specified internuclear distance $r$, is proportional to the projection $\Sigma$ of the total molecular spin on its axis. The split electronic term may therefore be written
\[
 U(r)+A(r)\Sigma,
\]

\SourcePage{391}{394}
where $U(r)$ is the energy of the original unsplit term and $A(r)$ is a function of $r$. This function depends on the original term, in particular on $\Lambda$, but not on $\Sigma$. Since the quantum number $\Omega$, rather than $\Sigma$, is normally used, it is more convenient to write $A\Omega$ in place of $A\Sigma$. The two expressions differ by $A\Lambda$, which may be included in $U(r)$. Thus the electronic term is
\begin{equation}U(r)+A(r)\Omega.\tag{83.3}\end{equation}

The components of the split term are equally spaced: neighboring components, whose $\Omega$ values differ by one, are separated by $A(r)$ independently of $\Omega$.

General considerations readily show that $A=0$ for $\Sigma$ terms. Apply time reversal. The energy must remain unchanged, while the directions of both the orbital and spin angular momenta about the axis are reversed. In the energy $A(r)\Sigma$, $\Sigma$ changes sign, so $A(r)$ must also change sign if the energy is to remain invariant. For $\Lambda\ne0$ this says nothing about the value of $A(r)$, since $A$ depends on the orbital angular momentum, which itself reverses. But when $\Lambda=0$, $A(r)$ must in any event remain unchanged and must therefore vanish identically. Hence, for $\Sigma$ terms, the spin--orbit interaction causes no first-order splitting. A splitting proportional to $\Sigma^2$ would appear only from the spin--orbit interaction in second order or the spin--spin interaction in first order, and is consequently small. This explains the fact noted above that $\Sigma$ terms generally belong to case $b$.

Once the multiplet splitting has been determined, molecular rotation can be included as a perturbation, in exact analogy with the derivation at the beginning of the preceding section. The rotational angular momentum of the nuclei is obtained by subtracting the electronic orbital angular momentum and spin from the total angular momentum. The centrifugal-energy operator is therefore
\[
 B(r)(\widehat{\mathbf J}-\widehat{\mathbf L}-\widehat{\mathbf S})^2.
\]
Averaging this quantity over the electronic state and adding it to (83.3), we obtain the required effective potential

\SourcePage{392}{395}
energy $U_J(r)$:
\[
\begin{split}
U_J(r)&=U(r)+A(r)\Omega+B(r)\overline{(\mathbf J-\mathbf L-\mathbf S)^2}\\
&=U(r)+A(r)\Omega+B(r)[\mathbf J^2-2\mathbf J(\overline{\mathbf L+\mathbf S})+\overline{\mathbf L^2}+2\overline{\mathbf{LS}}+\overline{\mathbf S^2}].
\end{split}
\]
The eigenvalue of $\mathbf J^2$ is $J(J+1)$. By the same reasoning as in \S~82,
\begin{equation}\overline{\mathbf L}=\mathbf n\Lambda,\qquad \overline{\mathbf S}=\mathbf n\Sigma,\tag{83.4}\end{equation}
and $(\widehat{\mathbf J}-\widehat{\mathbf L}-\widehat{\mathbf S})\mathbf n=0$, so that for the eigenvalues
\begin{equation}\mathbf J\mathbf n=(\mathbf L+\mathbf S)\mathbf n=\Lambda+\Sigma=\Omega.\tag{83.5}\end{equation}
Substitution gives
\[
U_J(r)=U(r)+A(r)\Omega+B(r)[J(J+1)-2\Omega^2+\overline{\mathbf L^2}+2\overline{\mathbf{LS}}+\overline{\mathbf S^2}].
\]
The average over the electronic state is taken with the zeroth-approximation wave functions\footnote{Zeroth order both in molecular rotation and in the spin--axis interaction.}. In this approximation spin is conserved, so $\mathbf S^2=S(S+1)$. The wave function is a product of a spin function and a coordinate function, and the angular momenta $\mathbf L$ and $\mathbf S$ are consequently averaged independently:
\[
\overline{\mathbf{LS}}=\Lambda\mathbf n\overline{\mathbf S}=\Lambda\Sigma.
\]
Finally, the mean square orbital angular momentum $\overline{\mathbf L^2}$ is independent of spin and is a function of $r$ characteristic of the given unsplit electronic term. All terms that depend on $r$ but not on $J$ and $\Sigma$ may be included in $U(r)$, while the term proportional to $\Sigma$, or equivalently to $\Omega$, may be absorbed into $A(r)\Omega$. The effective potential energy is therefore
\begin{equation}U_J(r)=U(r)+A(r)\Omega+B(r)[J(J+1)-2\Omega^2].\tag{83.6}\end{equation}

The molecular energy levels can now be obtained in the same way as in \S~82 from (82.5). Expanding $U(r)$ and $A(r)$ in powers of $\xi$, retaining terms through second order in $U(r)$ and only zeroth-order terms in the second and third terms, we find
\begin{equation}E=U_e+A_e\Omega+\hbar\omega_e\left(v+\frac12\right)+B_e[J(J+1)-2\Omega^2],\tag{83.7}\end{equation}

\SourcePage{393}{396}
where $A_e=A(r_e)$ and $B_e$ are constants characteristic of the given unsplit electronic term. Continuing the expansion would produce further terms of higher degree in the quantum numbers; we shall not write them here.
